\input{../universal-pre}
\newcommand{\mcal}[1]{\mathcal{#1}}
\title{\tbf{Quillen-Suslin Theorem}}
\author{\tbf{Notes} by Cong Wen, 06/2020}
\date{}

\begin{document}
\maketitle
\thispagestyle{empty}
\begin{abstract}
  This is my self-use notes\footnote{Last updated on June 27, 2020} on \tbf{Quillen-Suslin theorem}\cite{Quillen1976}, i.e., all the finite projective modules over polynomial rings are free. The proof is divided into two parts, the Hilbert-Serre theorem, which states that all the finite projective modules over polynomial rings are stably free, and the Quillen-Suslin theorem, all the stably free modules over polynomial rings are actually free.
\end{abstract}
\tableofcontents


% ----------------------------------------------------------------------------------------------------

\section{Algebraic Preliminaries}

This part collects some useful results in homological algebra and commutative algebra, and all the rings occurred below are commutative rings with unity.

\subsection{Easy Examples}
Surely every free module is projective, the Quillen-Suslin theorem states that if $R$ is a polynomial ring over field, then finite projective modules are also free. There are two easy examples of this. When $R$ is a PID or a local ring, any projective module is free. We presented here the proof of these two theorems.

\begin{lem}\label{lem:localfree}
	$(R,\mfk{m})$ is a local ring, $M$ is a finite $R$-Mod, then the size of any minimal generating set of $M$ is equal to $\on{dim}M\otimes_R\kappa(\mfk{m})$, in which $\kappa(\mfk{m})=R/\mfk{m}$.
\end{lem}
\begin{proof}
	If $\{m_1,\cdots, m_n\}$ is a minimal generating set of $M$, we only need to show $\{\oL{m}_1,\cdots,\oL{m}_n\}$ are linearly independent in $M/\mfk{m}M$ as a $\kappa(\mfk{m})$-Mod. To do this, suppose they are linearly dependent, WLOG, suppose $\{\oL{m}_2,\cdots,\oL{m}_n\}$ generates $M/\mfk{m}M$, then $\sum_{k=2}^nRm_k+\mfk{m}M=M$, by Nakayama's lemma, $M=\sum_{k=2}^nRm_k$, contradicting the choice of $\{m_k\}$. So $\{\oL{m}_1,\cdots,\oL{m}_n\}$ is a basis of the $\kappa(\mfk{m})$-vector space $M/\mfk{m}=M\otimes\kappa(\mfk{m})$, hence $n=\on{dim}M\otimes_R\kappa(\mfk{m})$.
\end{proof}

\begin{thm}\label{thm:freenessOnLocal}
	If $(R,\mfk{m})$ is a local ring and $M$ is a finite projective $R$-Mod, then $M$ is free.
\end{thm}
\begin{proof}
	Since $M$ is finite, take a minimal generating set $\{m_1,\cdots,m_n\}$, then the homomorphism $f:R^n\rightarrow M$ is surjective, and $R^n=M\oplus\on{ker}f$, by tensoring $\kappa(\mfk{m})$ on both side, we get a $\kappa(\mfk{m})$-vector space isomorphisms $\kappa(\mfk{m})^n=M\otimes_R\kappa(\mfk{m}) \oplus \on{ker}f\otimes_R\kappa(\mfk{m})$, but the dimension of $M\otimes_R\kappa(\mfk{m})$ is already $n$, so $\on{ker}f\otimes_R\kappa(\mfk{m})=0$, by Nakayama's lemma, $\on{ker}f=0$. Hence $M$ is free.
\end{proof}
\begin{rmk}
	Actually, Kaplansky proved the theorem also holds when $M$ is not finitely generated.
\end{rmk}

\begin{lem}\label{lem:subPIDfree}
	If $R$ is a principal ideal domain and $M$ is a free module, then every submodule of $M$ is free.
\end{lem}
\begin{rmk}
	We omit the proof here, in the case when $M$ is of finite rank, it is a typical linear algebra proof, in the case when $M$ is of infinite rank, Zorn's lemma has come into play. See details in \cite{Lang2002}.
\end{rmk}

\begin{thm}\label{thm:freenessonPID}
	If $R$ is a principal ideal domain and $M$ is a finite projective $R$-Mod, then $M$ is free.
\end{thm}
\begin{proof}
	Since $M$ is finite generated and projective, there is a split short exact sequence,
	$$
	0\rightarrow K\rightarrow R^n\rightarrow M\rightarrow 0.
	$$
	So $R^n = M\oplus K$, by Lem. \ref{lem:subPIDfree}, $M$ is free.
	
	In fact we can see that when $M$ is not finitely generated, the statement still holds.
\end{proof}
\begin{rmk}
	This is the trivial case of Quillen-Suslin theorem when the polynomial ring has only one variable. In the general proof below, we will show this fact in another manner, see Lem. \ref{PID}.
\end{rmk}

\subsection{Commutative Algebra}
\begin{lem}\label{lem:lift}
	Let $f:R\rightarrow S$ be a surjective ring homomorphism, $\mfk{p}\in\on{Spec}S$, denote by $\mfk{p}^c\equiv f^{-1}(\mfk{p})\in\on{Spec}R$. If $N\subset M$ are two $S$-Mod, and $M/N\cong S/\mfk{p}$, then $_{R}M/_{R}N = R/\mfk{p}^c$ in which $_{R}M$ is the lift of $M$ via map $f$.
\end{lem}
\begin{proof}
	Note that the lift of $_{S}M/_{S}N$ is clearly $_{R}M/_{R}N$, so it suffices to show the lift of $S/\mfk{p}$ is $R/\mfk{p}^c$. Define a map
	\begin{equation*}
		\begin{split}
			\varphi:R/\mfk{p}^c&\rightarrow S/\mfk{p}\\
			r+\mfk{p}^c&\mapsto f(r)+\mfk{p},
		\end{split}
	\end{equation*}
	then $\varphi$ is clearly well-defined, injective and surjective.
\end{proof}

\begin{thm}\label{proj&lf}
	$M$ is a finite $R$-Mod. $M$ is projective iff $M$ is locally free.
\end{thm}
\begin{proof}
	$\Rightarrow$: If $M$ is finite projective, then there exist $N$ such that $M\oplus N=R^n$, if $\mfk{p}$ is any prime ideal, then $M_\mfk{p}\oplus N_\mfk{p}=R_\mfk{p}^n$, so now $M_\mfk{p}$ is a finite projective $R_\mfk{p}$-Mod, by Lem. \ref{lem:localfree}, $M_\mfk{p}$ is a free $R_\mfk{p}$-Mod.
	
	$\Leftarrow$: To show $M$ is projective, it suffices to prove $\on{Hom}_R(M,-)$ is an exact functor. Let $0\rightarrow N'\rightarrow N\rightarrow N''$ be a short exact sequence of $R$-Mod. To show
	$$
	0\rightarrow\on{Hom}_R(M,N')\rightarrow\on{Hom}_R(M,N)\rightarrow\on{Hom}_R(M,N'')\rightarrow 0
	$$
	is exact, it suffices to prove for any prime $\mfk{p}$ the following is exact,
	$$
	0\rightarrow\on{Hom}_R(M,N')_\mfk{p}\rightarrow\on{Hom}_R(M,N)_\mfk{p}\rightarrow\on{Hom}_R(M,N'')_\mfk{p}\rightarrow 0,
	$$
	since $M$ is finite, it is equivalent to prove the following is exact,
	$$
	0\rightarrow\on{Hom}_{R_\mfk{p}}(M_\mfk{p},N'_\mfk{p})\rightarrow\on{Hom}_{R_\mfk{p}}(M_\mfk{p},N_\mfk{p})\rightarrow\on{Hom}_{R_\mfk{p}}(M_\mfk{p},N''_\mfk{p})\rightarrow 0.
	$$
	This sequence is exact because $0\rightarrow N'_\mfk{p}\rightarrow N_\mfk{p}\rightarrow N''_\mfk{p}$ is exact, and $M_\mfk{p}$ is finite free, hence $\on{Hom}_{R_\mfk{p}}(M_\mfk{p},-)$ is an exact functor.
	
	So $\on{Hom}_R(M,-)$ is an exact functor and $M$ is projective.	
\end{proof}

\begin{thm}
	$M$ is a finite $R$-Mod, then $\on{Supp}(M)=V(\on{ann}(M))$.
\end{thm}
\begin{proof}
	$\subset$: If $\mfk{p}\not\supset\on{ann}(M)$, then $A\backslash\mfk{p}\cap\on{ann}(M)\neq\emptyset$, hence $M_\mfk{p}=0$.
	
	$\supset$: If $\mfk{p}\supset\on{ann}(M)$, then $\on{ann}(M_\mfk{p})=\on{ann}(M)_\mfk{p}\subset\mfk{p}R_\mfk{p}\neq R_\mfk{p}$, hence $M_\mfk{p}\neq 0$ (Note the first equality holds only when $M$ is finite, otherwise, it only holds that $\on{ann}(M_\mfk{p})\supset\on{ann}(M)_\mfk{p}$).
\end{proof}

\begin{thm}\label{minimal}
	Let $R$ be a Noetherian ring, and $M$ be a $R$-module, then the following sets are the same:
	\begin{enumerate}
		\item Any filtration of $M=M_0\supsetneqq\cdots\supsetneqq M_n=0$ with $M_{i-1}/M_{i}\cong R/\mfk{p}_i$, the minimal elements of $\{\mfk{p}_i\}, 1\leq i\leq n$,
		\item Minimal elements of $\on{Supp}(M)$,
		\item Minimal elements of $\on{Ass}(M)$.
	\end{enumerate}
\end{thm}
\begin{proof}
	1. and 2. are the same:	use the finite filtration, we know $\on{Supp}(M)=\bigcup\on{Supp}(R/\mfk{p}_i)=\bigcup V(\mfk{p}_i)$, so 1. and 2. are the same.
	
	2. and 3. are the same: first we have $\on{Ass}(M)\subset\on{Supp}(M)$, to see this, for any $\mfk{p}=\on{ann}(m)$, $m/1$ is a non zero element in $M_\mfk{p}$, so $\mfk{p}\in\on{Supp}(M)$.
	
	Then it suffices to prove any minimal prime in $\on{Supp}(M)$ is an associated prime: any minimal prime in $\on{Supp}(M)$ is now also minimal in $\on{Ass}(M)$, for any minimal prime $\mfk{p}\in\on{Ass}(M)$, there exist a minimal prime $\mfk{q}\in\on{Supp}(M)$ such that $\mfk{q}\subset\mfk{p}$, but $\mfk{q}$ is also an associated prime, hence $\mfk{p}=\mfk{q}$.
	
	Next we prove any minimal prime in $\on{Supp}(M)$ is an associated prime, take a minimal prime $\mfk{p}\in\on{Supp}(M)$, since 1. and 2. are the same, denote $k$ the maximal index such that $\mfk{p}_k=\mfk{p}$, take $m\in M_{k-1}, m\notin M_k$, consider $\mfk{q}=\on{ann}(m)$, any element which annihilates $m$ also annihilates $\bar{m}\in M_{k-1}/M_k$, and any element of $\mfk{p}_{k}\cdots \mfk{p}_n$ annihilates $m$ because $m\mfk{p}_{k}\cdots \mfk{p}_n\subset M_k\mfk{p}_{k+1}\cdots \mfk{p}_n\subset \cdots\subset M_{n-1}\mfk{p}_n=\{0\}$. So we have
	$$
	\mfk{p}_{k}\cdots \mfk{p}_n\subset\mfk{q}\subset\mfk{p}_k.
	$$
	By the choice of $\mfk{p}$, we know $\mfk{p}_i\not\subset\mfk{p}, k+1\leq i\leq n$, so $\mfk{p}_{k+1}\cdots \mfk{p}_n\not\subset\mfk{p}$, take $f\in\mfk{p}_{k+1}\cdots \mfk{p}_n\backslash\mfk{p}$, we claim $\on{ann}(fm)=\mfk{p}$: If $afm=0$, then $af\in\mfk{q}\subset\mfk{p}_k$, hence $a\in\mfk{p}$, and any $a\in\mfk{p}$, $af\in\mfk{p}_{k}\cdots \mfk{p}_n\subset\mfk{q}$. So $\mfk{p}$ is an associated prime.
\end{proof}

\subsection{Homological Algebra}
\begin{dfn}
	$M$ is a finite $R$-Mod, then it is called stably free of type $m$ if there exists a free $R$-Mod $F$ of rank $m<\infty$ such that $M\oplus F$ is free.
\end{dfn}
\begin{rmk}\label{rem:SFMod}
	It is required $m=\on{rk}F<\infty$ since any projective module would be stably free if $m=\infty$ (By Eilenberg). It is also required in the definition that $M$ is finitely generated, this is because if $M$ is not finitely generated, $M$ will actually be free (by Gabel).
\end{rmk}
\begin{proof}[proof of Rem. \ref{rem:SFMod}]
	If $m=\infty$, for any projective module $P$, it is a summand of a free module $F$, write $Q\oplus P = F$, then $P\oplus F^{\mbb{N}} = P\oplus Q\oplus P\oplus\cdots = F^{\mbb{N}}$.
	
	If $M$ is not finitely generated, then let $M\oplus R^m=F=R^I$ in which $I$ is an infinite index set (otherwise $M$ will be finitely generated), and write $M=\on{ker}(F\xrightarrow{f} R^m)$. Since $R^m$ is finitely generated, there exists a finite subset $I_0\subset I$ such that $f(F_0)=R^m$, in which $F_0=R^{I_0}$. So $F_0+M=F$. Let $N=F_0\cap M=\on{ker}f|_{F_0}$, so there are exact sequences:
	\begin{align*}
		0\rightarrow N\rightarrow &F_0\rightarrow R^m\rightarrow 0\\
		0\rightarrow N\rightarrow &M\rightarrow M/N\rightarrow 0.
	\end{align*}
	Now $M/N=M/(F_0\cap M) = (F_0+M)/F_0 = F/F_0=R^{I\backslash I_0}$, so the sequences above split. Since $I\backslash I_0$ is still infinite, we can write $R^{I\backslash I_0}=R^m\oplus F'$.
	
	So $M=N\oplus M/N=N\oplus R^m\oplus F'=F_0\oplus F'$ is free.
\end{proof}

\begin{prp}
	Every stably free module $M$ determines a right invertible $m\times n, n>m$ matrix $A$ such that $M$ is the solution space of $A$ and vice versa (the solution space of every right invertible $m\times n$ matrix is stably free).
\end{prp}
\begin{proof}
	On the one hand, if $M$ is a stably free module, by definition and the remark, we got a split short exact sequence,
	$$
	0\rightarrow R^m\xrightarrow{f} R^n\rightarrow M\rightarrow 0.
	$$
	Since the sequence splits, there exist $g:R^n\rightarrow R^m$ such that $gf=\on{id}_{R^m}$ and $M=\on{ker}g$. The matrix of $g$ is then $A$. Note that $A$ is right invertible is equivalent to $g$ is surjective.
	
	On the other hand, given a right invertible $m\times n$ matrix $A$, denote by $M$ the solution space of $A$, the sequence below splits hence $M$ is stably free.
	$$
	0\rightarrow M\rightarrow R^n\xrightarrow{A} R^m\rightarrow 0.
	$$
\end{proof}

\begin{prp}
	$M=\on{ker}(R^n\xrightarrow{g} R^m), n>m$, in which $g$ is surjective, is free iff $g$ can be lifted to an isomorphism $\tilde{g}:R^n\rightarrow R^r\oplus R^m$ for some $r$ such that $\pi_2\tilde{g}=g$, in which $\pi_2:R^r\oplus R^m\rightarrow R^m$ is the projection map on the second factor.
\end{prp}
\begin{proof}
	The sequence
	$$
	0\rightarrow M\xrightarrow{i} R^n\xrightarrow{g} R^m\rightarrow 0.
	$$
	splits and is exact, let $j:R^n\rightarrow M$ be the homomorphism such that $ji=\on{id}_M$, then
	$$
	R^n\xrightarrow{\mat{j\\ g}}M\oplus R^m
	$$
	is an isomorphism.
	
	On the one hand if $v:M\cong R^r$, then $\tilde{g}=\mat{vj\\ g}$ is the lift of $g$. On the other hand if the lift of $g$ exists, then $M=\on{ker}g\cong\on{ker}\pi_2=R^r$
\end{proof}

\begin{lem}
	$R$ a non zero commutative ring, then $R$ has invariant base number (IBN) property. To be precise, if $R^m\cong R^n$ as $R$ modules, then $m=n$.
\end{lem}
\begin{rmk}
	So the $r$ occurred in the previous proposition must be $n-m$.
\end{rmk}
\begin{proof}
	Take any maximal ideal $\mfk{m}$ of $R$. Tensoring with $k=R/\mfk{m}$ gives a $k$-Mod isomorphism $k^m\cong k^n$. Isomorphism vector bases has the same dimension hence $m=n$.
\end{proof}

\begin{cor}\label{cor:matrix}
	$R$ a commutative ring, $A$ a right invertible $m\times n$ matrix, and $M$ is the solution space of $A$ (hence $M$ is stably free). $M$ is free iff $A$ can be extended to a $n\times n$ invertible matrix by adding $n-m$ rows.
\end{cor}

\begin{thm}\label{thm:sf2cplt}
	Let $R$ be any ring, then the following are equivalent,
	\begin{enumerate}
		\item Any stably free $R$-Mod is free,
		\item Any stably free $R$-Mod of type $1$ is free,
		\item Any unimodular row $f=(f_1,\cdots,f_n)$ in $R^n$ can be extended to an invertible matrix in $\on{GL}_n(R)$ by adding $n-1$ rows. (Such a row is called completable). Unimodular means $\sum_iRf_i=R$.
	\end{enumerate}
\end{thm}
\begin{proof}
	$1.\Rightarrow 2.$: trivial.
	
	$2.\Rightarrow 1.$: use induction, suppose any stably free $R$-Mod of type $r$, $r<n$ is free. $M$ a stably free $R$-Mod of type $n$, then $M\oplus R^n$ is free, so $M\oplus R$ is a stably free module of type $n-1$, hence free. So $M$ is stably free of type $1$ and hence free.
	
	$2.\Leftrightarrow 3.$: by Cor. \ref{cor:matrix}.
\end{proof}

\begin{thm}\label{thm:sf&fr}
	If $P$ is a projective $R$-Mod. $P$ is stably free iff $P$ admits a finite free resolution.
\end{thm}
\begin{proof}
	One direction is trivial, for the other direction, we use induction on the length of finite free resolution of $P$. Take the the finite free resolution as below
	$$
	0\rightarrow F_n\rightarrow\cdots\rightarrow F_0\rightarrow M\rightarrow 0.
	$$
	
	If $n\leq 1$, $P$ is surely stably free.
	
	Suppose the theorem holds for $n<k$, we prove for k:
	$$
	0\rightarrow F_k\rightarrow\cdots\rightarrow F_0\rightarrow M\rightarrow 0.
	$$
	Let $M_1=\on{ker}(F_0\rightarrow M)$, then split the long sequence into two:
	\begin{align*}
		0\rightarrow F_k\rightarrow\cdots&\rightarrow F_0\rightarrow M_1\rightarrow 0\\
		0\rightarrow M_1\rightarrow&F_0\rightarrow M\rightarrow 0
	\end{align*}
	By the induction hypothesis, $M_1$ is stably free, hence $M_1\oplus F$ is free for some free module $F$. Then $F_0\oplus F=M\oplus M_1\oplus F$ is free, hence $M$ is stably free.
\end{proof}

\begin{thm}
	$R$ is a Noetherian, then for a short exact sequence of $R$-Mod,
	$$
	0\rightarrow M'\rightarrow M\rightarrow M''\rightarrow0,
	$$
	then that any two of $M',M,M''$ admits a finite resolution implies the other also does.
\end{thm}
\begin{rmk}
	The proof of this theorem is omitted for three reasons: The theorem is itself easy to intuitively make sense, and the proof is rather typical in homological algebra, and it takes a few pages.
\end{rmk}

\section{Motivation}
\subsection{Serre-Swan Theorem}
In this part to demonstrate the motivation for this theorem, we states a bunch of theorems in geometry without proof and preliminaries. The readers are assumed to be familiar with knowledge in topology, differentiable manifold and algebraic geometry.

The basic intuition is that finite projective module corresponds to a finite rank vector bundles. On a topological space $M$, and a real topological vector bundle $E\rightarrow M$, the global sections $\Gamma(E)$ can surely be viewed as a $C(M)$-Mod, in which $C(M)$ is the ring of continuous function on $M$. Of course on a trivial bundle $F$, $\Gamma(F)$ is a free $C(M)$-Mod. Since every projective module is a direct summand of a free module, we can see the analogy from the following theorem:
\begin{thm}
	For each vector bundle $E\rightarrow B$, with $B$ a compact Hausdorff space, then there exists a vector bundle $E'$ such that $E\oplus E'$ is a trivial bundle.
\end{thm}
So every vector bundle can give rise to a finite projective $C(M)$-Mod, the other direction also holds, this is proved in 1962\cite{Swan1962}, Swan established the one-one correspondence in general case.
\begin{thm}
	Let $X$ be compact Hausdorff, then a $C(X)$-module $P$ is isomorphic to a module of the form $\Gamma(E)$ iff $P$ is finitely generated and projective.
\end{thm}

In fact, it is Serre's work which motivated the proof of Swan, and gave a one-one correspondence between algebraic vector bundles on affine scheme and finite projective modules. This can be found in Serre's famous paper \textit{Faisceaux Algebriques Coherents}\cite{Serre1955}, also in which he proposed the Serre's conjecture, i.e. the main theorem we are going to show in this note.

\begin{thm}
	If $\mcal{F}$ is a coherent sheaf on a connected affine scheme $X=\on{Spec}A$, then the following are equivalent:
	\begin{enumerate}
		\item $\mcal{F}(X)$ is a finite projective $A$-Mod,
		\item $\mcal{F}(X)$ is isomorphic to the sheaf of germs of sections of an algebraic vector bundle of base $X$.
	\end{enumerate}
\end{thm}

This theorem is a direct implication from Thm. \ref{proj&lf} below after getting familiar with the language of scheme theory.


In fact the philosophy can be further generalized and expressed in locally ringed space, and was proved by Morye\cite{Morye2013}. Let $(X,\mcal{O}_X)$ be any ringed space, and $\on{Lfb}(X)$ denote the full subcategory of $\mcal{O}_X$-Mod consisting of locally free $\mcal{O}_X$-modules of bounded rank. And $\on{Fgp}(A)$ denote the category of finite projective $A$-modules in which $A=\Gamma(X,\mcal{O}_X)$, then the Serre-Swan theorem can be stated in terms of ringed space:
\begin{dfn}
	Serre-Swan theorem holds for $(X,\mcal{O}_X)$ if the functor 
	
	$$
	\Gamma(X,-):\on{Lfb}(X)\rightarrow\on{Fgp}(\Gamma(X,\mcal{O}_X))
	$$
	
	is well-defined and an equivalence of categories.
\end{dfn}
Then it is proved that
\begin{thm}
	Let $(X,\mcal{O}_X)$ be a locally ringed space such that each $\mcal{F}$ in $\on{Lfb}(X)$ is acyclic and generated by finite many global sections, then the Serre-Swan theorem holds for $(X,\mcal{O}_X)$.
\end{thm}

We restate some of the examples in which Serre-Swan's theorem holds:
\begin{enumerate}
	\item (Serre's theorem) $(X,\mcal{O}_X)$ is an affine scheme.
	\item $(X,\mcal{O}_X)$ a ringed space, $X$ is paracompact of bounded topological dimension, and $\mcal{O}_X$ is a fine sheaf.
		\begin{itemize}
			\item (Swan's theorem) $(X,\mcal{C}_X)$ in which $X$ is a paracompact topological space with bounded topological dimension, and $\mcal{C}_X$ the sheaf of continuous real-valued functions.
			\item $(X,\mcal{C}^\infty(\mbb{R}))$ in which $X$ is a real smooth manifold of bounded dimension, $\mcal{C}^\infty(\mbb{R})$ the sheaf of smooth real-valued functions on $X$.
		\end{itemize}
	\item $(X,\mcal{O}_X)$ a locally ringed space, in which $X$ is compact.
	\item $(X,\mcal{O}_X)$ a Stein space.
\end{enumerate}


\subsection{Serre's Conjecture}
After the correspondence it is natural to think of the most common example $\mbb{A}^r_k=\on{Spec}k[x_1,\cdots,x_r]$. A classical result in topological vector bundles is that every vector bundle on a contractible base space is trivial, hence $k^r$ has no non-trivial bundles, and it is expected $\mbb{A}^r_k$ has the same property. Serre wrote that

\textit{On ignore s'il existe des $A$-moduls projectif de type fini qui ne soient pas libres, ou, ce qui revient au même, s'il existe des espaces fibrés algébraiques à fibres vectorielles, de base $K^r$, et non triviaux.}

This is known as Serre's conjecture, and was proved by Quillen and Suslin in 1976, later Leonid Vaseršteĭn later gave a simpler and much shorter proof of the theorem which is collected in this note.

\section{Proof}
\subsection{Hilbert-Serre Theorem}
\begin{thm}[Hilbert-Serre]
	$k$ is a field and $M$ is a finite projective module over $k[x_1,\cdots,x_n]$, then $M$ is stably free.
\end{thm}
This theorem is by Serre\cite{Serre1958} shortly after his conjecture in 1958. To prove this, we need another lemma.
\begin{lem}\label{lem:PreHSthm}
	$R$ is Noetherian, that every finite $R$-Mod admits a finite free resolution implies that every finite $R[x]$-Mod admits a finite free resolution.
\end{lem}
To see how this lemma implies Hilbert-Serre theorem, it suffices to notice every finite $k$-Mod is free, hence admits a finite free resolution, then every finite $k[x_1,\cdots,x_n]$-Mod admits a finite free resolution by the lemma. In particular, every projective finite $k[x_1,\cdots,x_n]$-Mod admits finite free resolution and is stably free by Thm. \ref{thm:sf&fr}.

\begin{proof}[Proof of Lem. \ref{lem:PreHSthm}]
	Let $M$ be a finite $R[x]$-Mod, then $R[x]$ is Noetherian and $M$ has a finite filtration
	$$
	M=M_0\supsetneqq M_1\supsetneqq\cdots\supsetneqq M_n=0,
	$$
	in which $M_{k-1}/M_k\cong R[x]/P_k$ and $P_k\in\on{Spec}R[x]$, now we have $n$ short exact sequences of $R[x]$-Mod,
	\begin{equation*}
		\begin{split}
			0\rightarrow M_n\rightarrow &M_{n-1}\rightarrow R[x]/P_n\rightarrow 0\\
			&\vdots\\
			0\rightarrow M_1\rightarrow &M_{0}\rightarrow R[x]/P_1\rightarrow 0
		\end{split}
	\end{equation*}
	
	To show $M_0=M$ admits a finite resolution, it suffices to show every $R[x]/P_k$ admits a finite free resolution, notice that the following short exact sequence as $R[x]$-Mod,
	$$
	0\rightarrow P_k\rightarrow R[x]\rightarrow R[x]/P_k\rightarrow0
	$$
	it then suffices to show every $P\in\on{Spec}R[x]$ admits a finite resolution, suppose not, we proceed by contradiction:
	
	Let $\mcal{S}$ be all prime ideals which does not admit a finite resolution, if $\mcal{S}$ is not empty, we can take $P\in\mcal{S}$ such that $\mfk{p}=P\cap R$ is a maximal element in $\{P\cap R:P\in\mcal{S}\}$. Let $P'=P/\mfk{p}[x], R'=R/\mfk{p}$, then $P'$ is a $R'[x]$-Mod. Since $P\subset R[x]$, $R[x]$ is Noetherian, $P$ is a finite $R[x]$-Mod, hence $P'$ is finite $R'[x]$-Mod.
	
	Take $f_1,\cdots f_n\in R'[x]$ as the generators of $P'$, and also take $f$ to be the polynomial of minimal degree in $P'$, and do factorization in $\on{Frac}(R')[x]$:
	$$
	f_i=q_if+r_i, r_i,q_i\in \on{Frac}(R')[x]
	$$
	take the common denominator of all $r_i$ and $q_i$ to be $d'\in R'$, since $R'$ is an integral domain ($\mfk{p}$ is prime), we know $d'\neq 0$. Back to $R'[x]$, we thus have $d'f_i=d'q_if+d'r_i, \on{deg}{d'r_i}<\on{deg}{f}$, since the degree of $f$ is minimal, we have $d'f_i=d'q_if$, which means $d'P'\subset fR'[x]\equiv (f)$.
	
	Let $N'=P'/(f)$, this is a finite $R'[x]$-Mod, also it can be lifted to a $R[x]$-Mod, denote it by $N$. Now we have a finite filtration $\{N_k'\}$ of $N'$, lift each one to $R[x]$-Mod, we get a finite filtration $\{N_k\}$ of $N$.
	\begin{equation*}
		\begin{split}
			N'=N_0'\supsetneqq N_1'\supsetneqq\cdots\supsetneqq N_r'=0,\\
			N=N_0\supsetneqq N_1\supsetneqq\cdots\supsetneqq N_r=0,
		\end{split}
	\end{equation*}
	such that $N_{k-1}'/N_k'=R'[x]/Q_k', N_{k-1}/N_k=R[x]/Q_k, Q_k=Q_k'^c$, this can be done because of Lem. \ref{lem:lift}. So $Q_k = Q_k'^c$ shows that $Q_k\supset\mfk{p}[x]$, hence $Q_k\cap R\supset\mfk{p}$. In addition, $d'$ annihilates $N'$, so the representative $d$ of $d'$ in $R$ annihilates $N$. So $d\in\cap\on{Ass}_{R[x]}(N)=\cap Q_k$, the last equality holds by Thm. \ref{minimal}. Also $d'\neq 0$ implies $d\notin\mfk{p}$, so for all $k$, $Q_k\cap R\supsetneqq\mfk{p}$. By the choice of $\mfk{p}$, we know every $Q_k$ admits a finite free resolution, hence $N_{k-1}/N_k\cong R[x]/Q_k$ also does.
	
	Notice that these short exact sequences for $N$
	\begin{equation*}
		\begin{split}
			0\rightarrow N_r\rightarrow &N_{r-1}\rightarrow R[x]/Q_r\rightarrow 0\\
			&\cdots\\
			0\rightarrow N_1\rightarrow &N_{0}\rightarrow R[x]/Q_1\rightarrow 0
		\end{split}
	\end{equation*}
	so $N=N_0$ admits a finite free resolution.
	
	Further, since $\mfk{p}$ is finite $R$-Mod, it admits a finite free resolution by condition:
	$$
	0\rightarrow R^{k_m}\rightarrow\cdots\rightarrow R^{k_0}\rightarrow\mfk{p}\rightarrow 0
	$$
	Now $R[x]$ is a free $R$-Mod, hence flat, so by tensoring with $R[x]$, we know that $\mfk{p}\otimes R[x]$ admits a finite resolution. But $\mfk{p}\otimes R[x]\cong\mfk{p}[x]$, this is because
	\begin{equation*}
		\begin{split}
			0\rightarrow\mfk{p}\rightarrow &R\rightarrow R/\mfk{p}\rightarrow0\\
			0\rightarrow\mfk{p}\otimes R[x]\rightarrow &R\otimes R[x]\rightarrow R/\mfk{p}\otimes R[x]\rightarrow0\\
		\end{split}
	\end{equation*}
	is exact, hence $\mfk{p}[x]$ admits a finite free resolution.
	
	In addition
	$$
	0\rightarrow\mfk{p}[x]\rightarrow R[x]\rightarrow R'[x]\rightarrow0
	$$
	is exact as $R[x]$-Mod, so $R'[x]$ as well as $(f)\equiv fR'[x]\cong R'[x]$ admits a finite free resolution. The isomorphism holds because $R'[x]$ is an integral domain.
	
	At last, by definition of $N$ and $P$, the sequences
	\begin{equation*}
		\begin{split}
			0\rightarrow(f)\rightarrow P'\rightarrow N\rightarrow0\\
			0\rightarrow\mfk{p}[x]\rightarrow P\rightarrow P'\rightarrow0\\
		\end{split}
	\end{equation*}	
	are exact, $(f)$ and $N$ admits finite free resolution, so $P'$ does, hence $P$ does, a contradiction!
	
	Overall, every $P\in\on{Spec}R[x]$ admits a finite free resolution, hence $M$ does.
\end{proof}
	
\subsection{Quillen-Suslin Theorem}
Based on the Hilbert-Serre theorem, we only need to prove every stably free module over a polynomial ring is free. Further by Thm. \ref{thm:sf2cplt}, this is equivalent to prove every unimodular row $f=(f_1,\cdots,f_n)$ in $k[x_1,\cdots,x_r]^n$ is completable.

By the definition of completable row, it is easy to see $f$ is completable if and only if there exists $M\in\on{GL}_n(k[x_1,\cdots,x_r])$ such that $fM = (1,0,\cdots,0)$. For two rows $u,v\in R^n$, we write $u\sim v$ if there exists $G\in\on{GL}_n(R)$ such that $v=uG$.

We proceed in four steps:
\begin{lem}\label{PID}
	$k$ is a field, $f=(f_1,\cdots,f_n)\in k[x]^n$ is unimodular, then $f$ is completable.
\end{lem}
\begin{proof}
	 To show $f(x)\sim (1,0,\cdots,0)$, we give a constructive proof:
	 
	 First, multiplying an invertible matrix of $\on{GL}_n(k[x])$ on a matrix is equivalent to do elementary column transformations on this matrix (this is true in $\on{GL}_n(R)$ when $R$ is an Euclidean domain, but not true in PID).
	 
	 In particular, swapping two component of $f$, adding a multiple of a component of $f$ to another component, and multiplying a component by a constant in $k$ these three operations can all be written in the form of multiplying by a invertible matrix on the right and vice versa.
	 
	Now if $f(x)$ is unimodular, we have $\on{gcd}(f_1(x),\cdots,f_n(x))=1$, so this guarantees that we can do Euclidean algorithm repeatedly to reduce $f$ to $(1,0,\cdots,0)$.
	
	To illustrate the process, take $n=2$, and let $f_1 = qf_2+r$, then
	$$
	\mat{f_1 & f_2}\mat{1&0\\-q&1} = \mat{r & f_2}
	$$
	now $\on{deg}r<\on{deg}f_1$, repeating the process, $r$ finally becomes $1$ (because $f_1$ and $f_2$ coprime), then
	$$
	\mat{1,f_2'}\mat{1&-f_2'\\0&1} = \mat{1 & 0}.
	$$
	So we can see that $f=(f_1,f_2)\sim (1,0)$, to be precise, $fM=(1,0)$, in which
	$$
	M=\mat{1&0\\-q&1}\mat{1&-f_2'\\0&1}
	$$
	
	As for $n>=2$, the process is similar.
\end{proof}
\begin{rmk}
	In the following, to prove a row is completable, we only need to construct a sequence of row operations to make it into $(1,0,\cdots,0)$.
\end{rmk}

\begin{thm}[Horrock]\label{thm:local}
	$(R,\mfk{m})$ is a local ring, $f=(f_1,\cdots,f_n)\in R[x]^n$ is unimodular and has a component with leading coefficient $1$, then $f$ is completable.
\end{thm}
\begin{proof}
	For convenience, denote by $\on{lc}(h)$ the leading coefficient of a polynomial $h$.

	When $n=1$, the only component is invertible and has leading coefficient $1$, hence is $1$.
	
	When $n=2$, write $f=(f_1,f_2)$, assume $uf_1+vf_2=1$, then
	$$
	(f_1,f_2)\mat{u & -f_2\\v & f_1} = (1,0).
	$$
	
	Assume now $n\geq 3$, we proceed by induction, define
	$$
	d=\on{min}\{\on{deg}f_i\mid \on{lc}(f_i)=1, 1\leq i\leq n\},
	$$
	and suppose the theorem holds for $d<d_0$, we prove for $d_0$ below:
	
	WLOG we assume $\on{deg}f_1=d, \on{lc}(f_1)=1$, using Euclidean algorithm. Now that $f_1$ can be used as a divider in the Euclidean algorithm, we can assume $\on{deg}f_i\leq d_0-1, i\geq 2$ (otherwise write $f_i = q_if_1+r_i$, then subtract $q_if_1$ from $f_i$).
	
	Claim: $f_i, i\geq 2$ cannot have all their coefficients in $\mfk{m}$. Suppose not, since $f$ is unimodular, there exists $g_i$, $\sum_{i=1}^nf_ig_i=1$. Passing into $R/\mfk{m}[x]$, we have $\oL{f}_1\oL{g}_1=\oL{1}$, since $f_1$ has leading coefficient $1$, then all the coefficient except the constant term of $g_1$ must be in $\mfk{m}$ and hence $\oL{g}_1=\oL{c}$, contradicting $\oL{f}_1\oL{g}_1=\oL{1}$.
	
	So WLOG assume $f_2$ has at least one coefficient $b_i$ not in $\mfk{m}$, hence invertible, write
	\begin{align*}
		&f_1=x^{d_0}+a_1x^{d_0-1}+\cdots +a_{d_0},\\
		&f_2=b_0x^s+\cdots+b_s, b_i\notin\mfk{m}, s\leq d_0-1.
	\end{align*}
	
	The goal now is to find a combination of $f_1,f_2$ with leading coefficient $1$ and degree less than $d_0$. To do this, define an ideal
	$$
	\mfk{a}=\{\on{lc}(h)\mid h=g_1f_1+g_2f_2, \on{deg}h<d_0, g_1,g_2\in R[x]\},
	$$
	to show $\mfk{a}$ is the whole ring, only need to show $\mfk{a}$ contains all the coefficients of $f_2$. This is achievable by simply playing with the two polynomials, for example $\on{lc}\left(0f_1+x^{d_0-s-1}f_2\right)=b_0$. In fact, let $a_0=1$, in general if we define
	$$
	h=-(b_0x^{r-1}+\cdots+b_{r-1})f_1+(a_0x^r+\cdots+a_r)x^{d_0-s-1}f_2, r\leq s,
	$$
	then $\on{lc}(h)=b_r$. This can easily be checked by expanding the coefficients.
	
	So $f\sim (h,f_2,\cdots, f_n), h=f_1g_1+f_2g_2, \on{deg} h\leq d_0-1$, by induction, $f\sim (1,0,\cdots,0)$ is completable.
\end{proof}
\begin{rmk}
	This theorem in fact generalizes the preceding lemma, allowing the coefficient to be in a local ring rather than a field.
\end{rmk}

\begin{thm}\label{thm:bivar}
	$R$ is an integral domain, $S$ a multiplicative set, $f=(f_1,\cdots,f_n)\in R[x]^n$, if $f(x)\sim f(0)$ over $S^{-1}R[x]$, then there exists $c\in S$, $f(x+cy)\sim f(x)$ over $R[x,y]$.
\end{thm}
\begin{proof}
	We start from $f(x)\sim f(0)$, let $M(x)\in\on{GL}_n(S^{-1}R[x])$ such that $f(x)M(x)=f(0)$, so $f(x)M(x)M(x+y)^{-1}=f(x+y)$, write $G(x,y)=M(x)M(x+y)^{-1}$, we see that $G(x,y)$ is a matrix with entries in $S^{-1}R[x,y]$, in addition $G(x,0)=\on{id}$.
	
	So $G(x,y)-G(x,0)=yH(x,y)$ in which $H$ is also a matrix with entries in $S^{-1}R[x,y]$, there then exist a element $c\neq 0$ (because $R$ is an integral domain) such that $cH(x,y)$ has coefficients in $R[x,y]$. So $G(x,cy) = \on{id}+cyH(x,cy)$ has all its entries in $R[x,y]$, by definition of $G$, we know
	$$
	f(x)G(x,cy)=f(x+cy).
	$$
	
	Now it only suffices to show $G(x,cy)$ is invertible over $R[x,y]$. Since $M(x)\in\on{GL}_n(S^{-1}R[x])$, we know $\on{det}M(x)\in (S^{-1}R)^*$ is a constant, hence the determinant of $G(x,y)$ is identically $1$, hence invertible. Done.
\end{proof}

\begin{thm}\label{ID}
	$R$ is an integral domain, $f=(f_1,\cdots,f_n)\in R[x]^n$ is unimodular and has a component with leading coefficient $1$, then $f(x)\sim f(0)$ over $R[x]$.
\end{thm}
\begin{proof}
	Let $J$ be the set of $c\in R$ such that $f(x+cy)\sim f(x)$ over $R[x,y]$, then $J$ is an ideal:
	
	If $a,b\in J$, then $f(x+ay+by)\sim f(x+ay)$ over $R[x+ay, y]$ and hence over $R[x,y]$ (Note $\on{GL}_n(R[x+ay,y])\subset\on{GL}_n(R[x,y])$). Further $f(x+ay)\sim f(x)$ over $R[x,y]$, hence $f(x+ay+by)\sim f(x)$ over $R[x,y]$ and $a+b\in J$.
	
	If $a\in J, r\in R$, then $f(x+ary)\sim f(x)$ over $R[x,ry]$ and hence over $R[x,y]$ (Note $\on{GL}_n(R[x,ry])\subset\on{GL}_n(R[x,y])$). So $ar\in J$.
	
	Claim: $J=R$. By Thm. \ref{thm:local}, $f(x)\sim f(0)$ over $R_\mfk{p}[x]$ for any prime $\mfk{p}$, followed by Thm. \ref{thm:bivar}, we know $f(x)\sim f(x+cy)$ over $R[x,y]$ in which $c\notin \mfk{p}$. Since $c\in J$, $J\not\subset\mfk{p}$. So $J$ is not contained in any maximal ideal and is hence the whole ring, so $1\in J$.
	
	Now $f(x+y)M(x,y)=f(x)$, in which $M(x,y)\in\on{GL}_n(R[x,y])$. Note that $M(0,y)$ is the image of $M(x,y)$ under the ring homomorphism $\tilde{\pi}:\on{GL}_n(R[x,y])\rightarrow\on{GL}_n(R[y])$ induced by $\pi:R[x,y]\rightarrow R[y]$, hence $M(0,y)$ is invertible.
	
	So by plugging into $x=0$ we get $f(y)\sim f(0)$.
\end{proof}

With the preceding theorems above, we can use induction to prove Quillen-Suslin theorem:
\begin{thm}[Quillen-Suslin]
	If $f=(f_1,\cdots,f_n)\in k[x_1,\cdots,x_r]^n$ is unimodular, then $f$ is completable.
\end{thm}
\begin{proof}
	Use induction on $r$:

	If $r=1$, by Lem. \ref{PID}.
	
	If for $r-1$ the theorem already holds, now prove for $r$. Take
	$$
	f_1(x_1,\cdots,x_r)=\sum_{(a)\in I} c_{a}x_1^{a_1}\cdots x_r^{a_r}, 0\neq c_a\in k
	$$
	we now need a $k$-algebra automorphism to make leading coefficient of $f_1$ (with respect to $x_r$) become $1$. To do this, take $d>\on{max}_{(a)\in I}\{a_i\}$, and let
	$$
	y_r = x_r, y_{r-1} = x_{r-1}-x_r^d, y_{r-2}=x_{r-2}-x_r^{d^2},\cdots, y_1=x_1-x_r^{d^{r-1}}.
	$$
	Substitute them in to $f_1$, we get
	$$
	f_1(y_1,\cdots,y_r)=\sum_{(a)\in I}c_ay_r^{a_r+a_{r-1}d+\cdots+a_1d^{r-1}}+g(y_1,\cdots,y_r),
	$$
	in which $g$ does not contain any pure power of $y_r$. Due to the choice of $d$, the degree in the first sum are all distinct, hence non zero. So now the leading coefficient of $f_1$ with respect to $y_r$ is in $k$ and hence invertible. (Note this is the same substitution trick in Noether Normalization Theorem)
	
	Now let $R=k[y_1\cdots,y_{r-1}]$, then every $f_i$ can be seen as a polynomial in $R[y_r]$, and $f_1$ has leading coefficient $1$. So by Thm \ref{ID}, $f=(f_1,\cdots, f_n)\sim f(0)$ over $R[y_r]$, i.e. there exists $M\in\on{GL}_n(R[y_r])=\on{GL}_n(k[y_1,\cdots,y_r])\subset\on{GL}_n(k[x_1,\cdots,x_r])$ such that $fM=f(0)\in R=k[y_1,\cdots,y_{r-1}]$.
	
	But $f(0)$ is still unimodular and is a polynomial in $r-1$ variables $\{y_1,\cdots,y_{r-1}\}$, so by the induction hypothesis, there exists $N\in\on{GL}_n(k[y_1,\cdots,y_{r-1}])\subset\on{GL}_n(k[x_1,\cdots,x_r])$, such that $f(0)N=(1,0,\cdots,0)$.
	
	Hence $f(x)MN = (1,0,\cdots,0)$, in which $MN\in \on{GL}_n(k[x_1,\cdots,x_r])$, done!
\end{proof}


\section*{Acknowledgement}
  I'd like to thank Prof. Pu Zhang and Prof. Mounir Hajli for their wonderful lectures on algebra, who encouraged and helped me a lot so that I had the confidence and knowledge to read through the whole proof.

% ----------------------------------------------------------------------------------------------------
\bibliographystyle{alpha}
\bibliography{../universal-ref}
\end{document}


